% Hindawi 期刊通用 LaTeX 模板（以 Shock and Vibration 为例）
% Hindawi 投稿系统接受标准 article class，基本要求：
%   - 单栏 (one column)
%   - 12pt 字体
%   - 双倍行距 (setspace)
%   - 1 英寸页边距 (geometry)
%   - 行号可选，部分期刊初审建议添加
% 编译方式：pdfLaTeX 或 XeLaTeX（中文作者可用 ctex 宏包）
\documentclass[12pt,a4paper]{article}

% ==================== 页面设置 ====================
\usepackage[margin=1in]{geometry}
\usepackage{setspace}
\doublespacing

% ==================== 基础宏包 ====================
\usepackage{amsmath,amssymb,amsfonts}
\usepackage{graphicx}
\usepackage{booktabs}
\usepackage{caption}
\usepackage{subcaption}
\usepackage{hyperref}
\usepackage{lineno}          % 行号，部分期刊初审建议保留
\linenumbers

% ==================== 标题与作者 ====================
\title{Vision-Based Motor Shaft Vibration Measurement Using STARK Tracking and Shaft-Diameter Self-Calibration}

\author{First Author$^{1,*}$, Second Author$^{2}$, Third Author$^{1}$ \\
\small $^{1}$School of Mechanical Engineering, University Name, City, Country \\
\small $^{2}$Department of Electrical Engineering, University Name, City, Country \\
\small $^{*}$Corresponding author: email@university.edu.cn}

\date{}

% ==================== 正文开始 ====================
\begin{document}

\maketitle

\begin{abstract}
\noindent
This study presents a non-contact vision-based method for measuring motor shaft vibration amplitudes.
A spatio-temporal transformer tracker (STARK-ST) is employed to localize the rotating shaft region,
followed by ellipse-aware circle detection for sub-pixel center estimation.
Pixel displacements are converted into physical coordinates through a shaft-diameter self-calibration strategy,
eliminating the need for a precise working distance.
Vibration amplitudes in the X and Y directions are computed as peak-to-peak values,
and a laser displacement sensor is used for cross-validation.
Experimental results demonstrate that the proposed system achieves comparable accuracy to the laser reference
while providing a low-cost and flexible alternative for rotating machinery monitoring.
\end{abstract}

\vspace{0.5em}
\noindent\textbf{Keywords:} visual tracking; motor shaft vibration; STARK; camera calibration; non-contact measurement

\section{Introduction}
Rotating machinery fault diagnosis relies heavily on vibration measurement.
Traditional contact sensors require mounting fixtures and cabling, which limit applicability in confined spaces.
Vision-based methods offer non-contact, full-field, and easy-deployment advantages.
However, high-speed rotation, motion blur, and illumination changes challenge robust target tracking.

In this work, we integrate the STARK-ST visual tracker with geometric shaft detection
and propose a shaft-diameter-based scale calibration method.
The contributions are as follows:
\begin{itemize}
    \item Application of a transformer-based tracker to rotating shaft vibration measurement.
    \item A robust center-detection pipeline combining STARK ROI prediction and ellipse/circle fitting.
    \item A scale self-calibration method using the known shaft diameter instead of a fixed working distance.
\end{itemize}

\section{Methodology}
\subsection{Visual Tracking}
The STARK-ST tracker predicts the target bounding box in each frame.
The predicted ROI is then refined by contour-based circle detection.

\subsection{Coordinate Conversion}
Given the detected pixel radius $r$ and the known shaft diameter $D$,
the scale factor $s$ is computed as
\begin{equation}
s = \frac{D}{2r} \quad \text{(m/pixel)}.
\end{equation}
The physical displacement relative to the optical center is
\begin{equation}
\Delta X = s \,(u - c_x), \qquad \Delta Y = s \,(v - c_y),
\end{equation}
where $(u,v)$ is the detected center and $(c_x,c_y)$ is the principal point.

\subsection{Vibration Amplitude}
The peak-to-peak amplitude is
\begin{equation}
A_x = \max(\Delta X) - \min(\Delta X), \qquad
A_y = \max(\Delta Y) - \min(\Delta Y),
\end{equation}
and the total amplitude is $A_{total} = \sqrt{A_x^2 + A_y^2}$.

\section{Experiments and Results}
Describe the experimental setup, dataset, comparison with laser sensor, and quantitative results.
Use tables and figures as needed.

\section{Conclusion}
Summarize the main findings and potential future work.

% ==================== 参考文献 ====================
\begin{thebibliography}{9}
\bibitem{stark2021}
B. Yan et al., ``Learning spatio-temporal transformer for visual tracking,'' \textit{Proc. IEEE/CVF ICCV}, 2021, pp. 10448--10457.

\bibitem{zhang2000}
Z. Zhang, ``A flexible new technique for camera calibration,'' \textit{IEEE Trans. Pattern Anal. Mach. Intell.}, vol. 22, no. 11, pp. 1330--1334, 2000.

\bibitem{visionvibration2020}
A. Author et al., ``Vision-based vibration measurement: A review,'' \textit{Mech. Syst. Signal Process.}, vol. 135, p. 106382, 2020.
\end{thebibliography}

\end{document}
